% LaTeX Canvas: Chapter Outline
\chapter{Remote Computing}
\label{ch:remote-computing}

\section*{Purpose}
Remote computing lets you use resources that are not physically on your laptop: university servers, high-performance clusters, lab machines, and cloud instances. This chapter teaches novices how to connect safely and reliably using SSH and VPNs, move files, run remote jobs, and use tunneling to access remote services (including Jupyter) as if they were local.

\section*{Learning objectives}
By the end of this chapter, you should be able to:
\begin{enumerate}
\item Explain client/server basics and what ``remote computing'' means.
\item Use SSH to connect to a remote machine and run commands interactively.
\item Use key-based authentication safely (generate keys, protect private keys, manage known hosts).
\item Transfer files with \texttt{scp} or \texttt{sftp} (and understand when to use GUI tools).
\item Use SSH tunneling (local, remote, and dynamic forwarding) for safe access to remote services.
\item Understand what a VPN does, when it is required, and common pitfalls.
\item Describe cloud computing at a practical level (instances, regions, networking rules, costs).
\item Apply baseline security practices: least privilege, MFA, key hygiene, and avoiding data leaks.
\end{enumerate}

\section*{Running theme: remote access is powerful, so make it boring}
The goal is not cleverness; it is predictable workflows that minimize risk.

\section{A beginner mental model}
\subsection{Local vs remote}
\begin{itemize}
\item \textbf{Local machine:} your laptop/desktop.
\item \textbf{Remote machine:} a computer you access over the network.
\item \textbf{Network:} the path between them (campus, home Wi-Fi, internet).
\end{itemize}

\subsection{Client/server language}
\begin{itemize}
\item \textbf{Client:} the tool you run locally (SSH client, VPN client, browser).
\item \textbf{Server:} the remote service or machine you connect to (SSH server, web server).
\end{itemize}

\subsection{Sessions, processes, and persistence}
\begin{itemize}
\item A remote \textbf{SSH session} is a connection + a shell.
\item Programs you start remotely may stop when your session ends unless you use job control tools.
\item Remote computers have their own file system, users, permissions, and installed software.
\end{itemize}

\subsection{Why remote computing exists}
\begin{itemize}
\item Access to more CPU/RAM/GPU.
\item Access to shared datasets and licensed tools.
\item Reliability (servers run longer than laptops).
\item Collaboration (shared environments).
\end{itemize}

\section{Prerequisites and safe defaults}
\subsection{Accounts and credentials}
\begin{itemize}
\item You need: a username, a host address, and an authentication method.
\item Prefer: institutional identity + MFA, and key-based SSH.
\end{itemize}

\subsection{Install/locate your tools}
\begin{itemize}
\item \textbf{macOS:} Terminal includes \texttt{ssh}, \texttt{scp}, \texttt{sftp}.
\item \textbf{Windows:} OpenSSH client is commonly available; otherwise enable/install it.
\item Optional: a code editor that supports remote development.
\end{itemize}

\subsection{Golden safety rules}
\begin{itemize}
\item Never share passwords or private keys.
\item Never paste secrets into commands or notebooks.
\item Assume your command history is readable by future you (and sometimes by administrators).
\item Use the minimum access needed; log out when done.
\end{itemize}

\section{SSH fundamentals}
\subsection{What SSH provides}
\begin{itemize}
\item Encrypted remote login.
\item Secure file transfer.
\item Secure tunneling of other network traffic.
\end{itemize}

\subsection{The simplest connection}
\begin{verbatim}
ssh username@hostname
\end{verbatim}
\begin{itemize}
\item You may need \texttt{-p <port>} if SSH uses a nonstandard port.
\item The first time you connect, SSH will ask about the host fingerprint.
\end{itemize}

\subsection{Host identity and \texttt{known\_hosts}}
\begin{itemize}
\item SSH remembers server fingerprints to prevent ``man-in-the-middle'' attacks.
\item Warning: a changed fingerprint may be benign (server rebuilt) or dangerous (impersonation).
\item Best practice: verify changes via a trusted channel (course staff / sysadmin).
\end{itemize}

\subsection{Authentication: passwords vs keys}
\begin{itemize}
\item \textbf{Password:} easy to start, less secure, often disabled on managed servers.
\item \textbf{Keys:} recommended; a private key stays on your device and proves identity.
\end{itemize}

\subsection{Key-based authentication basics}
\begin{itemize}
\item A key pair: \textbf{private key} (secret) + \textbf{public key} (shareable).
\item The public key is added to the server's authorized keys list.
\item Protect the private key with a passphrase.
\end{itemize}

\subsection{Generating a key pair (conceptual)}
\begin{verbatim}
ssh-keygen
\end{verbatim}
\begin{itemize}
\item Choose a strong passphrase.
\item Store keys in the default \texttt{~/.ssh/} location.
\item Keep your private key file permissions restrictive.
\end{itemize}

\subsection{SSH agent (optional but useful)}
\begin{itemize}
\item An agent holds decrypted keys in memory so you do not type your passphrase repeatedly.
\item Treat agent forwarding as advanced; do not enable it by default.
\end{itemize}

\subsection{Quality-of-life: \texttt{~/.ssh/config}}
\begin{itemize}
\item Create short nicknames for hosts.
\item Store per-host usernames, ports, and identity files.
\item Keep configuration readable; document nonstandard choices.
\end{itemize}

\section{File transfer: getting data in and out}
\subsection{The basic question: where is the file?}
\begin{itemize}
\item `Local path'' (your computer) vs `remote path'' (server).
\item Use absolute paths if you are unsure.
\end{itemize}

\subsection{Secure copy (\texttt{scp})}
\begin{itemize}
\item Good for quick transfers and scripts.
\item Learn directory copies (recursive) and quoting paths with spaces.
\end{itemize}

\subsection{SFTP (\texttt{sftp})}
\begin{itemize}
\item Interactive file transfer session.
\item Many GUI tools use SFTP under the hood.
\end{itemize}

\subsection{What not to do}
\begin{itemize}
\item Do not email datasets or copy them into unapproved cloud storage.
\item Do not move large files through a slow interactive session if a better transfer method exists.
\end{itemize}

\section{Tunneling and port forwarding (making remote services usable)}
\subsection{Why tunneling exists}
\begin{itemize}
\item Many services (Jupyter, databases, dashboards) bind to \texttt{localhost} on the server for safety.
\item A tunnel lets your local machine reach that service through an encrypted channel.
\end{itemize}

\subsection{Local forwarding (most common)}
\begin{itemize}
\item Scenario: remote server runs a service on \texttt{127.0.0.1:8888}; you want to open it in your local browser.
\end{itemize}
\begin{verbatim}
ssh -L 8888:localhost:8888 username@hostname
\end{verbatim}
\begin{itemize}
\item Then visit \texttt{[http://localhost:8888}](http://localhost:8888) in your local browser.
\item If \texttt{8888} is taken locally, choose another local port (e.g., 8899).
\end{itemize}

\subsection{Remote forwarding (advanced)}
\begin{itemize}
\item Scenario: you need the remote machine to reach a service running locally.
\item Use cautiously; it changes the exposure model.
\end{itemize}

\subsection{Dynamic forwarding (SOCKS proxy)}
\begin{itemize}
\item Scenario: route a browser or tool through the remote network.
\item Use only when you understand policies and implications.
\end{itemize}

\subsection{Tunnel hygiene}
\begin{itemize}
\item Use \texttt{-N} when you only need forwarding (no remote shell).
\item Prefer binding to \texttt{localhost} on your machine to avoid exposing ports to others.
\item Close tunnels when you are done.
\end{itemize}

\section{VPN fundamentals}
\subsection{What a VPN does (student-relevant)}
\begin{itemize}
\item Creates a secure network path to a private network (campus/lab).
\item Makes your device behave as if it is ``on campus'' for protected resources.
\item Often required to reach SSH hosts that are not publicly reachable.
\end{itemize}

\subsection{When you need a VPN}
\begin{itemize}
\item Accessing university resources from off campus.
\item Using internal-only hosts, file shares, or licensed services.
\item Connecting to a protected cluster login node.
\end{itemize}

\subsection{Common VPN pitfalls}
\begin{itemize}
\item VPN drops cause SSH disconnections.
\item Split-tunneling vs full-tunneling changes which traffic goes through the VPN.
\item Some VPNs interfere with local networking (printers, local dev servers).
\end{itemize}

\subsection{Practical advice}
\begin{itemize}
\item Connect VPN first, then SSH.
\item If you lose connection, reconnect VPN and re-establish SSH; do not assume remote processes survived.
\item Learn where your VPN client shows connection status and logs.
\end{itemize}

\section{Remote work patterns: from simple to professional}
\subsection{Pattern 1: login node + compute node (HPC concept)}
\begin{itemize}
\item You SSH into a \textbf{login node} for editing and job submission.
\item Actual computation runs on separate nodes via a scheduler.
\item Rule: do not run heavy jobs on login nodes.
\end{itemize}

\subsection{Pattern 2: remote development}
\begin{itemize}
\item Edit locally and run remotely, or edit remotely using a remote extension.
\item Keep projects in version control; avoid manual copying back and forth.
\end{itemize}

\subsection{Pattern 3: notebooks on servers}
\begin{itemize}
\item Start Jupyter remotely.
\item Use SSH local port forwarding to access it.
\item Prefer passwords/tokens and binding to \texttt{localhost}.
\end{itemize}

\section{Cloud computing basics (what novices must know)}
\subsection{What ``the cloud'' is}
\begin{itemize}
\item Renting computing resources (virtual machines, storage, managed services).
\item You control enough to break things, including security.
\end{itemize}

\subsection{Core terms}
\begin{description}
\item[Instance/VM] A rented computer.
\item[Region/zone] Where the machine physically resides.
\item[Storage] Object storage (buckets) vs attached disks.
\item[Networking rules] Firewalls/security groups controlling inbound/outbound.
\end{description}

\subsection{A safe ``first cloud VM'' workflow}
\begin{enumerate}
\item Create the instance in a region.
\item Create or choose an SSH key pair.
\item Configure inbound access: allow SSH only from your IP (not from the whole internet).
\item Connect via SSH.
\item Shut down/terminate when finished to control cost.
\end{enumerate}

\subsection{Cost and risk}
\begin{itemize}
\item Cloud resources cost money while running (and sometimes while stored).
\item Set budget alerts if available.
\item Do not expose services publicly unless required and reviewed.
\end{itemize}

\section{Security best practices for remote computing}
\subsection{Key hygiene}
\begin{itemize}
\item Use passphrases; do not store private keys in shared folders.
\item Avoid key sprawl: use separate keys per purpose if required, and retire unused keys.
\item If a key might be compromised, rotate it immediately.
\end{itemize}

\subsection{Least privilege and separation}
\begin{itemize}
\item Use non-admin accounts for daily work.
\item Use sudo/admin only when necessary.
\item Keep data access minimal and policy-compliant.
\end{itemize}

\subsection{Safe networking defaults}
\begin{itemize}
\item Restrict inbound ports (SSH only; no public databases).
\item Prefer tunnels over opening ports.
\item Keep software patched on remote machines.
\end{itemize}

\subsection{Logging and traceability (student level)}
\begin{itemize}
\item Keep a small ``remote session log'': what host, what you did, what changed.
\item Use version control for code changes.
\end{itemize}

\section{Troubleshooting and diagnostics}
\subsection{Connection failures}
\begin{itemize}
\item Check: internet, VPN status, host name, username.
\item Check: port and firewall rules.
\item Check: key permissions and correct key selection.
\end{itemize}

\subsection{Host key warnings}
\begin{itemize}
\item Do not ignore.
\item Verify with a trusted source before accepting changes.
\end{itemize}

\subsection{Timeouts and unstable connections}
\begin{itemize}
\item Use keep-alives if permitted.
\item Prefer running long tasks via schedulers or detached sessions.
\end{itemize}

\subsection{``Permission denied''}
\begin{itemize}
\item It may be the wrong username.
\item It may be the wrong key.
\item It may be a server-side access control issue.
\end{itemize}

\section{Worked examples (outline)}
\subsection{Example 1: First SSH login}
\begin{itemize}
\item Verify prerequisites (account, host, VPN).
\item Connect and interpret the login banner.
\item Navigate remote directories and run a simple command.
\end{itemize}

\subsection{Example 2: Move a dataset to a server}
\begin{itemize}
\item Identify local and remote paths.
\item Transfer with \texttt{scp} or \texttt{sftp}.
\item Verify integrity (file size, checksums if taught).
\end{itemize}

\subsection{Example 3: Remote Jupyter via SSH tunnel}
\begin{itemize}
\item Start Jupyter on the server bound to \texttt{localhost}.
\item Create a local tunnel.
\item Connect in browser; close tunnel when done.
\end{itemize}

\subsection{Example 4: Launch and secure a small cloud VM}
\begin{itemize}
\item Create instance; restrict inbound SSH to your IP.
\item SSH in; run a short job.
\item Terminate instance; confirm costs stop accruing.
\end{itemize}

\section{Exercises}
\begin{enumerate}
\item Generate an SSH key pair with a passphrase and identify where keys are stored.
\item Create an \texttt{~/.ssh/config} entry for a host and connect using the alias.
\item Transfer a directory of files to a server and verify the transfer.
\item Create a local port forward and access a remote service in your browser.
\item Connect to a required campus resource through a VPN and note how your network behavior changes.
\item Cloud practice (if allowed): launch a VM, SSH in, then shut it down and confirm termination.
\end{enumerate}

\section{One-page checklist}
\begin{itemize}
\item I know whether I need a VPN before SSH.
\item I can state: username, hostname, and authentication method.
\item My private key is protected (passphrase, correct permissions, not shared).
\item I understand host key prompts and do not ignore warnings.
\item I can transfer files securely and verify destination paths.
\item I can set up a local SSH tunnel for remote services.
\item I restrict network exposure (no unnecessary public ports).
\item I log out and terminate cloud resources when finished.
\end{itemize}

% \appendix
\section{Quick reference: common SSH commands}
\begin{verbatim}
ssh username@hostname
ssh -p 2222 username@hostname
ssh -i ~/.ssh/my_key username@hostname
ssh -L 8888:localhost:8888 username@hostname
scp localfile username@hostname:/remote/path/
sftp username@hostname
\end{verbatim}

\section{Quick reference: vocabulary}
\begin{description}
\item[SSH] Secure Shell protocol for encrypted remote access.
\item[VPN] Virtual Private Network; encrypted connection to a private network.
\item[Tunnel/Port forward] Using SSH to carry traffic for another service.
\item[Public key] Shareable key used by the server to verify the corresponding private key.
\item[Private key] Secret key that proves your identity; must be protected.
\item[Host key] Server identity fingerprint stored in \texttt{known_hosts}.
\item[Security group/firewall] Rules controlling allowed inbound/outbound network traffic.
\end{description}
